% Fragmento público generado desde propietarios íntegros.
% Residencia: 52.9.1.
% Fuente: ../incoming/radion_monografia_20260827/manuscrito/sections/10_minkowski_hodge_lorentz.tex:56-100
\subsubsection{Dualidad de Hodge y operador de cuarto de giro}

Fijadas la firma \((-+++ )\) y una orientación, la estrella de Hodge sobre
dos-formas satisface
\begin{equation}
\star:\Lambda^2Q_4^*\to\Lambda^2Q_4^*,
\qquad \star^2=-I.
\label{eq:hodge-star}
\end{equation}
De este modo, el cuarto de giro que en el plano de fase se expresaba mediante
\(J^2=-I\) reaparece en el espacio de campos como estructura compleja sobre
\(\Lambda^2\). Los dominios son diferentes, pero la genealogía de orientación
se conserva mediante el funtor.

Sea \(A_{\mathrm{em}}\) un potencial y
\begin{equation}
F=\diff A_{\mathrm{em}}.
\label{eq:F-dA}
\end{equation}
La descomposición de \(F\) respecto a un observador temporal produce campos
eléctrico y magnético; \(\star F\) intercambia sus papeles con el signo fijado
por la orientación y la firma. La polaridad constitutiva de
\eqref{eq:EM-conjugation} adquiere así una realización diferencial.

\paragraph{Cadena de transducción electromagnética y causal}

El transporte que debe mantenerse visible es
\begin{equation}
\Gamma_9
\longrightarrow\text{hélice de memoria}
\longrightarrow(Q_4,B_{\mathrm{caus}},\star)
\longrightarrow A_{\mathrm{em}}
\longrightarrow F=\diff A_{\mathrm{em}}
\longrightarrow
m\nabla_u u=q(\iota_uF)^\sharp.
\label{eq:lorentz-chain}
\end{equation}
Cada flecha añade estructura: la holonomía conserva retorno; el cociente fija
causalidad; Hodge fija dualidad; el potencial fija el campo; el acoplamiento
con carga produce dinámica.

La ecuación final es la fuerza de Lorentz en forma geométrica. La hélice TPK
es su antecedente de memoria, no una órbita física ya movida por un campo. La
dinámica aparece sólo al introducir \(F\), \(q\) y la conexión \(\nabla\).

